\section{Four lessons networking already paid for}

Resource management under the tension between per-unit guarantees and
aggregate efficiency is the central design problem of packet networking's
middle decades, and the field's resolutions were expensive: standardized
architectures were built, deployed at scale, and abandoned before the
surviving doctrine settled. Four of those lessons transfer to the serving
layer, and the discipline of Section 4 is their deliberate transfer. We
state each in its original setting, because the failure modes that
produced them are already visible, in miniature, in the systems of
Section 2.

Lesson one: per-flow guarantees do not scale to all flows, so guarantee
few and aggregate the rest. The Integrated Services architecture offered
end-to-end per-flow reservations, signaled by RSVP, to any flow that
asked \citep{rfc2205}; it failed in deployment substantially because per-flow
state in every element of the path could not scale. Its successor
inverted the design: Differentiated Services keeps per-flow complexity at
the edge and gives the core only a small set of aggregate classes
\citep{rfc2475}. The surviving doctrine is neither ``no guarantees'' nor
``guarantees for all''; it is guarantees rationed to the few flows whose
requirements justify per-flow state, with class-based service for the
remainder. Section 4 adopts this shape directly: a deliberately bounded
reserved subset, and aggregate classes for everything else.

Lesson two: load shedding is a design surface, and it should act early,
probabilistically, and per class. Random Early Detection drops packets
with a probability that rises as smoothed queue occupancy grows, before
the queue is full, precisely so that overload is signaled gradually
rather than discovered at exhaustion \citep{floyd1993red}; its
class-differentiated descendants shed lower-precedence traffic first
within the same discipline \citep{rfc2597}. The alternative, tail drop at a
full buffer, produces synchronized collapse and serves the most recent
arrival as badly as the most important one. Section 2's schedulers shed
load today, but deterministically and at the cliff: demotion at a
deadline, rejection at a ceiling, preemption at exhaustion. Section 4
transfers the early-probabilistic-per-class shape, with one deliberate
change of signal: the serving layer possesses a quantity networks lack, a
ledger of forward commitments, and Section 4 keys the shedding
probability to committed headroom rather than to observed queue state.

Lesson three: the plane that supervises a system must not share fate with
the system it supervises. Operators reach a wedged or compromised machine
through a management path that is powered, connected, and authenticated
independently of the machine's own stack: the baseboard management
controller and its descendants \citep{ipmi_spec}. The design premise is
unconditional: the moment supervision runs inside the thing supervised,
every failure of the thing is potentially a failure of its supervision,
and the cases where supervision matters most are exactly the cases where
the system is least trustworthy. Today's model-level oversight, guard
models in the pipeline, self-critique in the context, probes in the
process, sits inside the supervised stack. Section 4 transfers the
doctrine: a management plane whose authority and availability are
independent of the inference infrastructure, whatever signals it consumes
from within.

Lesson four: bind authority to provenance at the edge, not to content
downstream. Networks learned to discard traffic whose claimed origin does
not match the path it arrived by, at the ingress edge, before that
traffic can spend the network's resources or authority \citep{rfc2827,rfc3704}. The alternative, judging packets by their payloads deeper in
the network, loses to any adversary who can forge plausible content,
which is every adversary. The serving-layer analog is stated carefully in
Section 4, because content-level defenses for language models are an
active field with real results; the transfer claims only the resource
half: whatever a flow's content says, the provenance class assigned at
admission should bound what capacity and authority the flow can hold.

One table summarizes the mapping (Appendix C expands it): rationed
per-flow guarantees over aggregate classes $\to$ the bounded reserved subset;
early probabilistic class-aware shedding $\to$ headroom-keyed degradation;
fate-independent supervision $\to$ the management plane; provenance-bound
authority at the edge $\to$ trust-keyed admission. The transfers hold at the
granularity where they are claimed, the flow, where a control decision is
cheap relative to the minutes or hours of work it governs; they fail at
per-token or per-layer granularity, where the decision rivals the work;
and flows are not packets, since a flow's accumulated state entangles it
with the machinery serving it in ways a datagram never is. The lessons
are a design inheritance, bought at someone else's expense, and the
discipline that follows spends it. The evaluation that follows the
discipline then tests whether the inheritance survives the transfer: the
first lesson's economics, this paper finds, largely do not (Sections 7
and 9), the second lesson's transferred signal proves untenable in its
transferred form, and the third and fourth lessons' properties were never
economic claims and are examined on their own terms in Section 9.
